% Options for packages loaded elsewhere
\PassOptionsToPackage{unicode}{hyperref}
\PassOptionsToPackage{hyphens}{url}
\documentclass[
]{article}
\usepackage{xcolor}
\usepackage[margin=1in]{geometry}
\usepackage{amsmath,amssymb}
\setcounter{secnumdepth}{-\maxdimen} % remove section numbering
\usepackage{iftex}
\ifPDFTeX
  \usepackage[T1]{fontenc}
  \usepackage[utf8]{inputenc}
  \usepackage{textcomp} % provide euro and other symbols
\else % if luatex or xetex
  \usepackage{unicode-math} % this also loads fontspec
  \defaultfontfeatures{Scale=MatchLowercase}
  \defaultfontfeatures[\rmfamily]{Ligatures=TeX,Scale=1}
\fi
\usepackage{lmodern}
\ifPDFTeX\else
  % xetex/luatex font selection
\fi
% Use upquote if available, for straight quotes in verbatim environments
\IfFileExists{upquote.sty}{\usepackage{upquote}}{}
\IfFileExists{microtype.sty}{% use microtype if available
  \usepackage[]{microtype}
  \UseMicrotypeSet[protrusion]{basicmath} % disable protrusion for tt fonts
}{}
\makeatletter
\@ifundefined{KOMAClassName}{% if non-KOMA class
  \IfFileExists{parskip.sty}{%
    \usepackage{parskip}
  }{% else
    \setlength{\parindent}{0pt}
    \setlength{\parskip}{6pt plus 2pt minus 1pt}}
}{% if KOMA class
  \KOMAoptions{parskip=half}}
\makeatother
\usepackage{color}
\usepackage{fancyvrb}
\newcommand{\VerbBar}{|}
\newcommand{\VERB}{\Verb[commandchars=\\\{\}]}
\DefineVerbatimEnvironment{Highlighting}{Verbatim}{commandchars=\\\{\}}
% Add ',fontsize=\small' for more characters per line
\usepackage{framed}
\definecolor{shadecolor}{RGB}{248,248,248}
\newenvironment{Shaded}{\begin{snugshade}}{\end{snugshade}}
\newcommand{\AlertTok}[1]{\textcolor[rgb]{0.94,0.16,0.16}{#1}}
\newcommand{\AnnotationTok}[1]{\textcolor[rgb]{0.56,0.35,0.01}{\textbf{\textit{#1}}}}
\newcommand{\AttributeTok}[1]{\textcolor[rgb]{0.13,0.29,0.53}{#1}}
\newcommand{\BaseNTok}[1]{\textcolor[rgb]{0.00,0.00,0.81}{#1}}
\newcommand{\BuiltInTok}[1]{#1}
\newcommand{\CharTok}[1]{\textcolor[rgb]{0.31,0.60,0.02}{#1}}
\newcommand{\CommentTok}[1]{\textcolor[rgb]{0.56,0.35,0.01}{\textit{#1}}}
\newcommand{\CommentVarTok}[1]{\textcolor[rgb]{0.56,0.35,0.01}{\textbf{\textit{#1}}}}
\newcommand{\ConstantTok}[1]{\textcolor[rgb]{0.56,0.35,0.01}{#1}}
\newcommand{\ControlFlowTok}[1]{\textcolor[rgb]{0.13,0.29,0.53}{\textbf{#1}}}
\newcommand{\DataTypeTok}[1]{\textcolor[rgb]{0.13,0.29,0.53}{#1}}
\newcommand{\DecValTok}[1]{\textcolor[rgb]{0.00,0.00,0.81}{#1}}
\newcommand{\DocumentationTok}[1]{\textcolor[rgb]{0.56,0.35,0.01}{\textbf{\textit{#1}}}}
\newcommand{\ErrorTok}[1]{\textcolor[rgb]{0.64,0.00,0.00}{\textbf{#1}}}
\newcommand{\ExtensionTok}[1]{#1}
\newcommand{\FloatTok}[1]{\textcolor[rgb]{0.00,0.00,0.81}{#1}}
\newcommand{\FunctionTok}[1]{\textcolor[rgb]{0.13,0.29,0.53}{\textbf{#1}}}
\newcommand{\ImportTok}[1]{#1}
\newcommand{\InformationTok}[1]{\textcolor[rgb]{0.56,0.35,0.01}{\textbf{\textit{#1}}}}
\newcommand{\KeywordTok}[1]{\textcolor[rgb]{0.13,0.29,0.53}{\textbf{#1}}}
\newcommand{\NormalTok}[1]{#1}
\newcommand{\OperatorTok}[1]{\textcolor[rgb]{0.81,0.36,0.00}{\textbf{#1}}}
\newcommand{\OtherTok}[1]{\textcolor[rgb]{0.56,0.35,0.01}{#1}}
\newcommand{\PreprocessorTok}[1]{\textcolor[rgb]{0.56,0.35,0.01}{\textit{#1}}}
\newcommand{\RegionMarkerTok}[1]{#1}
\newcommand{\SpecialCharTok}[1]{\textcolor[rgb]{0.81,0.36,0.00}{\textbf{#1}}}
\newcommand{\SpecialStringTok}[1]{\textcolor[rgb]{0.31,0.60,0.02}{#1}}
\newcommand{\StringTok}[1]{\textcolor[rgb]{0.31,0.60,0.02}{#1}}
\newcommand{\VariableTok}[1]{\textcolor[rgb]{0.00,0.00,0.00}{#1}}
\newcommand{\VerbatimStringTok}[1]{\textcolor[rgb]{0.31,0.60,0.02}{#1}}
\newcommand{\WarningTok}[1]{\textcolor[rgb]{0.56,0.35,0.01}{\textbf{\textit{#1}}}}
\usepackage{graphicx}
\makeatletter
\newsavebox\pandoc@box
\newcommand*\pandocbounded[1]{% scales image to fit in text height/width
  \sbox\pandoc@box{#1}%
  \Gscale@div\@tempa{\textheight}{\dimexpr\ht\pandoc@box+\dp\pandoc@box\relax}%
  \Gscale@div\@tempb{\linewidth}{\wd\pandoc@box}%
  \ifdim\@tempb\p@<\@tempa\p@\let\@tempa\@tempb\fi% select the smaller of both
  \ifdim\@tempa\p@<\p@\scalebox{\@tempa}{\usebox\pandoc@box}%
  \else\usebox{\pandoc@box}%
  \fi%
}
% Set default figure placement to htbp
\def\fps@figure{htbp}
\makeatother
\ifLuaTeX
  \usepackage{luacolor}
  \usepackage[soul]{lua-ul}
\else
  \usepackage{soul}
\fi
\setlength{\emergencystretch}{3em} % prevent overfull lines
\providecommand{\tightlist}{%
  \setlength{\itemsep}{0pt}\setlength{\parskip}{0pt}}
\usepackage{bookmark}
\IfFileExists{xurl.sty}{\usepackage{xurl}}{} % add URL line breaks if available
\urlstyle{same}
\hypersetup{
  pdftitle={TD\_Devoir\_Simulation\_Aléatoire},
  hidelinks,
  pdfcreator={LaTeX via pandoc}}

\title{TD\_Devoir\_Simulation\_Aléatoire}
\author{}
\date{\vspace{-2.5em}}

\begin{document}
\maketitle

\subsubsection{\texorpdfstring{\ul{EXERCICE1}}{EXERCICE1}}\label{exercice1}

\paragraph{\texorpdfstring{\textbf{1.Donner le code pour acquérir les
données et les
visualiser.}}{1.Donner le code pour acquérir les données et les visualiser.}}\label{donner-le-code-pour-acquuxe9rir-les-donnuxe9es-et-les-visualiser.}

\begin{Shaded}
\begin{Highlighting}[]
\CommentTok{\#1{-} Code pour l\textquotesingle{}aquisition des données}

\NormalTok{Obs}\OtherTok{=}\FunctionTok{as.vector}\NormalTok{(LakeHuron)}



\FunctionTok{plot}\NormalTok{(Obs, }\AttributeTok{type=}\StringTok{"l"}\NormalTok{, }\AttributeTok{main=}\StringTok{"Niveaux d\textquotesingle{}eau LakeHuron (1875{-}1972)"}\NormalTok{)}
\end{Highlighting}
\end{Shaded}

\pandocbounded{\includegraphics[keepaspectratio]{TP_SIMULATION_files/figure-latex/unnamed-chunk-1-1.pdf}}

\paragraph{\texorpdfstring{\textbf{2. Tracer la fonction de répartition
empirique.}}{2. Tracer la fonction de répartition empirique.}}\label{tracer-la-fonction-de-ruxe9partition-empirique.}

\begin{Shaded}
\begin{Highlighting}[]
\CommentTok{\#2 Tracons la fonction empirique}
\NormalTok{Fonction\_empirique }\OtherTok{\textless{}{-}} \FunctionTok{ecdf}\NormalTok{(Obs) }
 \CommentTok{\# Tracé }
\FunctionTok{plot}\NormalTok{(Fonction\_empirique, }\AttributeTok{main =} \StringTok{"Fonction de répartition empirique (ECDF)"}\NormalTok{,}\AttributeTok{xlab =} \StringTok{"x"}\NormalTok{, }\AttributeTok{ylab =} \StringTok{"F(x)"}\NormalTok{, }\AttributeTok{col =} \StringTok{"blue"}\NormalTok{, }\AttributeTok{lwd =} \DecValTok{2}\NormalTok{) }

\FunctionTok{grid}\NormalTok{()}
\end{Highlighting}
\end{Shaded}

\pandocbounded{\includegraphics[keepaspectratio]{TP_SIMULATION_files/figure-latex/unnamed-chunk-2-1.pdf}}

\paragraph{\texorpdfstring{3\textbf{. \ul{Ajouter sur ce graphique la
fonction de répartition de la loi normale dont l'espérance et
l'écart-type sont la moyenne et l'écart-type de la variable
Obs.}}}{3. Ajouter sur ce graphique la fonction de répartition de la loi normale dont l'espérance et l'écart-type sont la moyenne et l'écart-type de la variable Obs.}}\label{ajouter-sur-ce-graphique-la-fonction-de-ruxe9partition-de-la-loi-normale-dont-lespuxe9rance-et-luxe9cart-type-sont-la-moyenne-et-luxe9cart-type-de-la-variable-obs.}

\begin{Shaded}
\begin{Highlighting}[]
\CommentTok{\#3 Ajoute la fonction de répartition de la loi normale de moyenne et d\textquotesingle{}ecart type celle de Obs}


\CommentTok{\#Moyenne et ecart type}

\NormalTok{mu}\OtherTok{=}\FunctionTok{mean}\NormalTok{(Obs)}

\NormalTok{ecart\_type}\OtherTok{=}\FunctionTok{sd}\NormalTok{(Obs)}


\NormalTok{y}\OtherTok{=}\FunctionTok{pnorm}\NormalTok{(Obs,mu,ecart\_type)}

 

\FunctionTok{plot}\NormalTok{(Fonction\_empirique, }\AttributeTok{main =} \StringTok{"Fonction de répartition empirique (ECDF)"}\NormalTok{, }
\AttributeTok{xlab =} \StringTok{"x"}\NormalTok{, }\AttributeTok{ylab =} \StringTok{"F(x)"}\NormalTok{, }\AttributeTok{col =} \StringTok{"blue"}\NormalTok{, }\AttributeTok{lwd =} \DecValTok{2}\NormalTok{) }

\FunctionTok{curve}\NormalTok{(}\FunctionTok{pnorm}\NormalTok{(x, mu, }\AttributeTok{sd=}\NormalTok{ecart\_type), }\AttributeTok{col=}\StringTok{"red"}\NormalTok{, }\AttributeTok{lwd=}\DecValTok{2}\NormalTok{, }\AttributeTok{add=}\ConstantTok{TRUE}\NormalTok{)}

\FunctionTok{legend}\NormalTok{(}\StringTok{"topleft"}\NormalTok{, }\AttributeTok{legend=}\FunctionTok{c}\NormalTok{(}\StringTok{"Empirique"}\NormalTok{, }\StringTok{"Normale"}\NormalTok{),}
       \AttributeTok{col=}\FunctionTok{c}\NormalTok{(}\StringTok{"blue"}\NormalTok{, }\StringTok{"red"}\NormalTok{), }\AttributeTok{lwd=}\DecValTok{2}\NormalTok{)}

\CommentTok{\#FONCTION DE REPARTITION}



\FunctionTok{grid}\NormalTok{()}
\end{Highlighting}
\end{Shaded}

\pandocbounded{\includegraphics[keepaspectratio]{TP_SIMULATION_files/figure-latex/unnamed-chunk-3-1.pdf}}

\paragraph{\texorpdfstring{\textbf{4. Tracer la fonction de quantile
empirique.}}{4. Tracer la fonction de quantile empirique.}}\label{tracer-la-fonction-de-quantile-empirique.}

\begin{Shaded}
\begin{Highlighting}[]
\CommentTok{\#Fonction de quantile}


\NormalTok{N}\OtherTok{=}\FunctionTok{length}\NormalTok{(}\FunctionTok{sort}\NormalTok{(Obs))}

\NormalTok{ptx }\OtherTok{=} \FunctionTok{seq}\NormalTok{(}\AttributeTok{from=}\DecValTok{1}\SpecialCharTok{/}\NormalTok{N, }\AttributeTok{to=}\DecValTok{1}\NormalTok{,}\AttributeTok{by=}\DecValTok{1}\SpecialCharTok{/}\NormalTok{N) }

\NormalTok{Qn }\OtherTok{=} \FunctionTok{quantile}\NormalTok{(Obs, ptx) }

\FunctionTok{plot}\NormalTok{(ptx, Qn, }\AttributeTok{type=}\StringTok{\textquotesingle{}l\textquotesingle{}}\NormalTok{, }\AttributeTok{xlab=}\StringTok{"p"}\NormalTok{,}\AttributeTok{ylab=}\StringTok{"Qn(p)"}\NormalTok{, }\AttributeTok{main=}\StringTok{"Fonction de quantile empirique"}\NormalTok{, }\AttributeTok{col=}\DecValTok{4}\NormalTok{, }
\AttributeTok{lwd=} \DecValTok{2}\NormalTok{)}
\end{Highlighting}
\end{Shaded}

\pandocbounded{\includegraphics[keepaspectratio]{TP_SIMULATION_files/figure-latex/unnamed-chunk-4-1.pdf}}

\paragraph{\texorpdfstring{\textsuperscript{\ul{\textbf{5-Tracer~le~QQ-plot~de~la~variable~Obs.}}}}{5-Tracer~le~QQ-plot~de~la~variable~Obs.}}\label{tracer-le-qq-plot-de-la-variable-obs.}

\begin{Shaded}
\begin{Highlighting}[]
\CommentTok{\#QQ{-}plot}
\FunctionTok{qqnorm}\NormalTok{(Obs, }\AttributeTok{xlab=}\StringTok{"Q{-}theo"}\NormalTok{, }\AttributeTok{ylab=}\StringTok{"Q{-}emp"}\NormalTok{, }\AttributeTok{col=}\DecValTok{4}\NormalTok{, }\AttributeTok{main=}\StringTok{"QQ{-}plot"}\NormalTok{)}
\end{Highlighting}
\end{Shaded}

\pandocbounded{\includegraphics[keepaspectratio]{TP_SIMULATION_files/figure-latex/unnamed-chunk-5-1.pdf}}

\paragraph{\texorpdfstring{\textsuperscript{\ul{\textbf{6-Tracer~l'histogramme~de~la~variable~Obs~ainsi~que~la~densité~de~la~loi~normale~dont~l'espérance~et~l'écart-type~sont~la~moyenne~et~l'écart-type~de~la~variable~Obs.}}}}{6-Tracer~l'histogramme~de~la~variable~Obs~ainsi~que~la~densité~de~la~loi~normale~dont~l'espérance~et~l'écart-type~sont~la~moyenne~et~l'écart-type~de~la~variable~Obs.}}\label{tracer-lhistogramme-de-la-variable-obs-ainsi-que-la-densituxe9-de-la-loi-normale-dont-lespuxe9rance-et-luxe9cart-type-sont-la-moyenne-et-luxe9cart-type-de-la-variable-obs.}

\begin{Shaded}
\begin{Highlighting}[]
\FunctionTok{hist}\NormalTok{(Obs, }\AttributeTok{prob=}\NormalTok{T,}\AttributeTok{col=}\StringTok{"lightgray"}\NormalTok{,}\AttributeTok{main=}\StringTok{"Histogramme + densité normale"}\NormalTok{)}
\CommentTok{\#Densité normale}

\FunctionTok{curve}\NormalTok{(}\FunctionTok{dnorm}\NormalTok{(x, mu, ecart\_type), }\AttributeTok{col=}\StringTok{"red"}\NormalTok{, }\AttributeTok{lwd=}\DecValTok{2}\NormalTok{, }\AttributeTok{add=}\ConstantTok{TRUE}\NormalTok{)}
\end{Highlighting}
\end{Shaded}

\pandocbounded{\includegraphics[keepaspectratio]{TP_SIMULATION_files/figure-latex/unnamed-chunk-6-1.pdf}}

\paragraph{\texorpdfstring{\textsuperscript{\textbf{7.Calculer~puis~représenter~un~estimateur~à~noyau~de~la~densité.}}}{7.Calculer~puis~représenter~un~estimateur~à~noyau~de~la~densité.}}\label{calculer-puis-repruxe9senter-un-estimateur-uxe0-noyau-de-la-densituxe9.}

\begin{Shaded}
\begin{Highlighting}[]
\NormalTok{dens }\OtherTok{\textless{}{-}} \FunctionTok{density}\NormalTok{(Obs)  }\CommentTok{\#Calcul}

\FunctionTok{plot}\NormalTok{(dens, }\AttributeTok{main=}\StringTok{"Estimateur à noyau"}\NormalTok{) }\CommentTok{\#Répresentation}
\end{Highlighting}
\end{Shaded}

\pandocbounded{\includegraphics[keepaspectratio]{TP_SIMULATION_files/figure-latex/unnamed-chunk-7-1.pdf}}

\paragraph{\texorpdfstring{\textsuperscript{\ul{\textbf{-Quel~est~le~noyau~utilisé~?}}}}{-Quel~est~le~noyau~utilisé~?}}\label{quel-est-le-noyau-utilisuxe9}

Le noyau de la densité utilisé est le noyau gaussien

\paragraph{\texorpdfstring{\textsuperscript{\ul{\textbf{-Quelle~est~la~fenêtre~utilisée~?}}}}{-Quelle~est~la~fenêtre~utilisée~?}}\label{quelle-est-la-fenuxeatre-utilisuxe9e}

\begin{Shaded}
\begin{Highlighting}[]
\CommentTok{\#fenetre utilise}
\NormalTok{dens}\SpecialCharTok{$}\NormalTok{bw}
\end{Highlighting}
\end{Shaded}

\begin{verbatim}
## [1] 0.4671343
\end{verbatim}

\textbf{La fenêtre utilisé est} :

\begin{verbatim}
0.4671343
\end{verbatim}

\paragraph{\texorpdfstring{\textsuperscript{\ul{\textbf{-Tracer~dans~une~même~fenêtre~plusieurs~estimateurs~à~noyau.}}}}{-Tracer~dans~une~même~fenêtre~plusieurs~estimateurs~à~noyau.}}\label{tracer-dans-une-muxeame-fenuxeatre-plusieurs-estimateurs-uxe0-noyau.}

\begin{Shaded}
\begin{Highlighting}[]
\NormalTok{dens\_gauss }\OtherTok{\textless{}{-}} \FunctionTok{density}\NormalTok{(Obs, }\AttributeTok{kernel=}\StringTok{"gaussian"}\NormalTok{)}
\NormalTok{dens\_rect  }\OtherTok{\textless{}{-}} \FunctionTok{density}\NormalTok{(Obs, }\AttributeTok{kernel=}\StringTok{"rectangular"}\NormalTok{)}
\NormalTok{dens\_tri   }\OtherTok{\textless{}{-}} \FunctionTok{density}\NormalTok{(Obs, }\AttributeTok{kernel=}\StringTok{"triangular"}\NormalTok{)}

\FunctionTok{plot}\NormalTok{(dens\_gauss, }\AttributeTok{col=}\StringTok{"blue"}\NormalTok{, }\AttributeTok{lwd=}\DecValTok{2}\NormalTok{, }\AttributeTok{main=}\StringTok{"Plusieurs estimateurs à noyau"}\NormalTok{)}

\FunctionTok{lines}\NormalTok{(dens\_rect,  }\AttributeTok{col=}\StringTok{"red"}\NormalTok{, }\AttributeTok{lwd=}\DecValTok{2}\NormalTok{)}

\FunctionTok{lines}\NormalTok{(dens\_tri,   }\AttributeTok{col=}\StringTok{"green"}\NormalTok{, }\AttributeTok{lwd=}\DecValTok{2}\NormalTok{)}

\FunctionTok{legend}\NormalTok{(}\StringTok{"topright"}\NormalTok{, }\AttributeTok{legend=}\FunctionTok{c}\NormalTok{(}\StringTok{"Gaussien"}\NormalTok{,}\StringTok{"Rectangulaire"}\NormalTok{,}\StringTok{"Triangulaire"}\NormalTok{),}
       \AttributeTok{col=}\FunctionTok{c}\NormalTok{(}\StringTok{"blue"}\NormalTok{,}\StringTok{"red"}\NormalTok{,}\StringTok{"green"}\NormalTok{), }\AttributeTok{lwd=}\DecValTok{2}\NormalTok{)}
\end{Highlighting}
\end{Shaded}

\pandocbounded{\includegraphics[keepaspectratio]{TP_SIMULATION_files/figure-latex/unnamed-chunk-9-1.pdf}}

\subsection{\texorpdfstring{\textsuperscript{EXERCICE2}}{EXERCICE2}}\label{exercice2}

\paragraph{\texorpdfstring{\ul{\textbf{Utilisons la méthode bootstrap
pour déterminer intervalle de confiance à 95 \% des moyennes
échantillonnales pour différentes valeurs de
B}}}{Utilisons la méthode bootstrap pour déterminer intervalle de confiance à 95 \% des moyennes échantillonnales pour différentes valeurs de B}}\label{utilisons-la-muxe9thode-bootstrap-pour-duxe9terminer-intervalle-de-confiance-uxe0-95-des-moyennes-uxe9chantillonnales-pour-diffuxe9rentes-valeurs-de-b}

\begin{Shaded}
\begin{Highlighting}[]
\FunctionTok{library}\NormalTok{(boot) }
\end{Highlighting}
\end{Shaded}

\begin{verbatim}
## Warning: le package 'boot' a été compilé avec la version R 4.5.2
\end{verbatim}

\begin{Shaded}
\begin{Highlighting}[]
\CommentTok{\# Données }
\NormalTok{x }\OtherTok{\textless{}{-}} \FunctionTok{c}\NormalTok{(}\DecValTok{51}\NormalTok{, }\DecValTok{45}\NormalTok{, }\DecValTok{49}\NormalTok{, }\DecValTok{66}\NormalTok{, }\DecValTok{53}\NormalTok{, }\DecValTok{41}\NormalTok{, }\DecValTok{58}\NormalTok{, }\DecValTok{56}\NormalTok{, }\DecValTok{60}\NormalTok{, }\DecValTok{63}\NormalTok{, }\DecValTok{75}\NormalTok{, }\DecValTok{89}\NormalTok{, }\DecValTok{73}\NormalTok{, }\DecValTok{84}\NormalTok{, }\DecValTok{66}\NormalTok{, }\DecValTok{85}\NormalTok{, }\DecValTok{73}\NormalTok{, }\DecValTok{71}\NormalTok{, }\DecValTok{78}\NormalTok{, }\DecValTok{65}\NormalTok{) }

\FunctionTok{set.seed}\NormalTok{(}\DecValTok{123}\NormalTok{)}
\NormalTok{fct }\OtherTok{\textless{}{-}} \ControlFlowTok{function}\NormalTok{(data, i) }\FunctionTok{mean}\NormalTok{(data[i]) }\CommentTok{\# Bootstrap}
\end{Highlighting}
\end{Shaded}

\begin{Shaded}
\begin{Highlighting}[]
 \CommentTok{\#  Intervalle par la Methode bootstrap pour B=50}

\NormalTok{b1 }\OtherTok{\textless{}{-}} \FunctionTok{boot}\NormalTok{(x, fct, }\AttributeTok{R =} \DecValTok{50}\NormalTok{)}
\FunctionTok{boot.ci}\NormalTok{(b1)}
\end{Highlighting}
\end{Shaded}

\begin{verbatim}
## Warning in boot.ci(b1): les variances de bootstrap sont nécessaires pour les
## intervalles studentisés
\end{verbatim}

\begin{verbatim}
## Warning in norm.inter(t, adj.alpha): statistiques d'ordre extrême utilisées
## comme points finaux
\end{verbatim}

\begin{verbatim}
## BOOTSTRAP CONFIDENCE INTERVAL CALCULATIONS
## Based on 50 bootstrap replicates
## 
## CALL : 
## boot.ci(boot.out = b1)
## 
## Intervals : 
## Level      Normal              Basic         
## 95%   (59.67, 70.20 )   (59.85, 72.39 )  
## 
## Level     Percentile            BCa          
## 95%   (57.71, 70.25 )   (56.80, 69.44 )  
## Calculations and Intervals on Original Scale
## Some basic intervals may be unstable
## Some percentile intervals may be unstable
## Warning : BCa Intervals used Extreme Quantiles
## Some BCa intervals may be unstable
\end{verbatim}

\begin{Shaded}
\begin{Highlighting}[]
 \CommentTok{\# Intervalle par la Methode bootstrap pour B=200}

\NormalTok{b2 }\OtherTok{\textless{}{-}} \FunctionTok{boot}\NormalTok{(x, fct, }\AttributeTok{R =} \DecValTok{200}\NormalTok{)}
\NormalTok{Intervalle\_confiance2 }\OtherTok{\textless{}{-}}\FunctionTok{boot.ci}\NormalTok{(b2)}
\end{Highlighting}
\end{Shaded}

\begin{verbatim}
## Warning in boot.ci(b2): les variances de bootstrap sont nécessaires pour les
## intervalles studentisés
\end{verbatim}

\begin{Shaded}
\begin{Highlighting}[]
\NormalTok{Intervalle\_confiance2}
\end{Highlighting}
\end{Shaded}

\begin{verbatim}
## BOOTSTRAP CONFIDENCE INTERVAL CALCULATIONS
## Based on 200 bootstrap replicates
## 
## CALL : 
## boot.ci(boot.out = b2)
## 
## Intervals : 
## Level      Normal              Basic         
## 95%   (59.04, 70.69 )   (58.41, 70.24 )  
## 
## Level     Percentile            BCa          
## 95%   (59.86, 71.69 )   (59.66, 71.45 )  
## Calculations and Intervals on Original Scale
## Some basic intervals may be unstable
## Some percentile intervals may be unstable
## Some BCa intervals may be unstable
\end{verbatim}

\begin{Shaded}
\begin{Highlighting}[]
 \CommentTok{\# Intervalle par la Methode bootstrap pour B=1000}

\NormalTok{b3 }\OtherTok{\textless{}{-}} \FunctionTok{boot}\NormalTok{(x, fct, }\AttributeTok{R =} \DecValTok{1000}\NormalTok{)}
\NormalTok{Intervalle\_confiance3 }\OtherTok{\textless{}{-}}\FunctionTok{boot.ci}\NormalTok{(b3)}
\end{Highlighting}
\end{Shaded}

\begin{verbatim}
## Warning in boot.ci(b3): les variances de bootstrap sont nécessaires pour les
## intervalles studentisés
\end{verbatim}

\begin{Shaded}
\begin{Highlighting}[]
\NormalTok{Intervalle\_confiance3}
\end{Highlighting}
\end{Shaded}

\begin{verbatim}
## BOOTSTRAP CONFIDENCE INTERVAL CALCULATIONS
## Based on 1000 bootstrap replicates
## 
## CALL : 
## boot.ci(boot.out = b3)
## 
## Intervals : 
## Level      Normal              Basic         
## 95%   (59.40, 70.78 )   (59.85, 71.35 )  
## 
## Level     Percentile            BCa          
## 95%   (58.75, 70.25 )   (58.55, 70.10 )  
## Calculations and Intervals on Original Scale
\end{verbatim}

\begin{Shaded}
\begin{Highlighting}[]
 \CommentTok{\#Intervalle par la Methode bootstrap pour B=2000}

\NormalTok{b4 }\OtherTok{\textless{}{-}} \FunctionTok{boot}\NormalTok{(x, fct, }\AttributeTok{R =} \DecValTok{2000}\NormalTok{)}
\NormalTok{Intervalle\_confiance4 }\OtherTok{\textless{}{-}}\FunctionTok{boot.ci}\NormalTok{(b4)}
\end{Highlighting}
\end{Shaded}

\begin{verbatim}
## Warning in boot.ci(b4): les variances de bootstrap sont nécessaires pour les
## intervalles studentisés
\end{verbatim}

\begin{Shaded}
\begin{Highlighting}[]
\NormalTok{Intervalle\_confiance4}
\end{Highlighting}
\end{Shaded}

\begin{verbatim}
## BOOTSTRAP CONFIDENCE INTERVAL CALCULATIONS
## Based on 2000 bootstrap replicates
## 
## CALL : 
## boot.ci(boot.out = b4)
## 
## Intervals : 
## Level      Normal              Basic         
## 95%   (59.33, 71.02 )   (59.25, 71.15 )  
## 
## Level     Percentile            BCa          
## 95%   (58.95, 70.85 )   (59.15, 70.90 )  
## Calculations and Intervals on Original Scale
\end{verbatim}

\ul{\textbf{Conclusion:}}

Pour B=50

\ul{Intervalle de confiance}: \textbf{{[}59.67, 70.20 {]}}

Pour B=200

\ul{Intervalle de confiance}: \textbf{\textbar59.40, 71.03 {]}}

Pour B=1000

\ul{Intervalle de confiance}: \textbf{\textbar59.31, 71.05 {]}}

Pour B=2000

\ul{Intervalle de confiance}: \textbf{\textbar59.07, 71.01{]}}

\ul{\textbf{On réalise le constat suivant lequel l'interavlle de
confiance varie entre : {[}59.07; 71.20{]}}}

\begin{quote}
\end{quote}

\subsection{\texorpdfstring{\textsuperscript{\ul{EXERCICE3}}}{EXERCICE3}}\label{exercice3}

\subsubsection{\texorpdfstring{\textsuperscript{1-Tracer~la~fonction~de~répartition~empirique.}}{1-Tracer~la~fonction~de~répartition~empirique.}}\label{tracer-la-fonction-de-ruxe9partition-empirique.-1}

\begin{Shaded}
\begin{Highlighting}[]
\CommentTok{\#1{-} Code pour l\textquotesingle{}aquisition des données}

\NormalTok{data2}\OtherTok{=}\FunctionTok{as.vector}\NormalTok{(rivers)}
\NormalTok{Obs1}\OtherTok{=}\FunctionTok{sort}\NormalTok{(data2)}

\CommentTok{\#2 Tracons la fonction empirique}
\NormalTok{Fonction\_empirique2 }\OtherTok{\textless{}{-}} \FunctionTok{ecdf}\NormalTok{(Obs1) }
 \CommentTok{\# Tracé }
\FunctionTok{plot}\NormalTok{(Fonction\_empirique2, }\AttributeTok{main =} \StringTok{"Fonction de répartition empirique (ECDF)"}\NormalTok{, }
\AttributeTok{xlab =} \StringTok{"x"}\NormalTok{, }\AttributeTok{ylab =} \StringTok{"F(x)"}\NormalTok{, }\AttributeTok{col =} \StringTok{"blue"}\NormalTok{, }\AttributeTok{lwd =} \DecValTok{2}\NormalTok{) }

\FunctionTok{grid}\NormalTok{()}
\end{Highlighting}
\end{Shaded}

\pandocbounded{\includegraphics[keepaspectratio]{TP_SIMULATION_files/figure-latex/unnamed-chunk-15-1.pdf}}

\subsubsection{\texorpdfstring{\textsuperscript{2-Superposer~la~fonction~de~répartition~de~la~loi~normale.}}{2-Superposer~la~fonction~de~répartition~de~la~loi~normale.}}\label{superposer-la-fonction-de-ruxe9partition-de-la-loi-normale.}

\begin{Shaded}
\begin{Highlighting}[]
\CommentTok{\#2 Ajoute la fonction de répartition de la loi normale de moyenne et d\textquotesingle{}ecart type celle de Obs}


\CommentTok{\#Moyenne et ecart type}

\NormalTok{mu1}\OtherTok{=}\FunctionTok{mean}\NormalTok{(Obs1)}

\NormalTok{ecart\_type1}\OtherTok{=}\FunctionTok{sd}\NormalTok{(Obs1)}


\FunctionTok{plot}\NormalTok{(Fonction\_empirique2, }\AttributeTok{main =} \StringTok{"Fonction de répartition empirique (ECDF)"}\NormalTok{, }\AttributeTok{xlab =} \StringTok{"x"}\NormalTok{, }\AttributeTok{ylab =} \StringTok{"F(x)"}\NormalTok{, }\AttributeTok{col =} \StringTok{"blue"}\NormalTok{, }\AttributeTok{lwd =} \DecValTok{2}\NormalTok{) }


\CommentTok{\#FONCTION DE REPARTITION}


\FunctionTok{curve}\NormalTok{(}\FunctionTok{pnorm}\NormalTok{(x,}\AttributeTok{mean=}\NormalTok{mu1,}\AttributeTok{sd=}\NormalTok{ecart\_type1),}\AttributeTok{col =} \StringTok{"red"}\NormalTok{, }\AttributeTok{lwd =} \DecValTok{2}\NormalTok{,}\AttributeTok{add=}\ConstantTok{TRUE}\NormalTok{)}

\FunctionTok{legend}\NormalTok{(}\StringTok{"topleft"}\NormalTok{, }\AttributeTok{legend=}\FunctionTok{c}\NormalTok{(}\StringTok{"Empirique"}\NormalTok{, }\StringTok{"Normale"}\NormalTok{),}
       \AttributeTok{col=}\FunctionTok{c}\NormalTok{(}\StringTok{"blue"}\NormalTok{, }\StringTok{"red"}\NormalTok{), }\AttributeTok{lwd=}\DecValTok{2}\NormalTok{)}


\FunctionTok{grid}\NormalTok{()}
\end{Highlighting}
\end{Shaded}

\pandocbounded{\includegraphics[keepaspectratio]{TP_SIMULATION_files/figure-latex/unnamed-chunk-16-1.pdf}}

\subsubsection{\texorpdfstring{\textsuperscript{3-~Faire~un~test~d'ajustement~de~Kolmogorov-Smirnov~et~Shapiro-Wilks.}}{3-~Faire~un~test~d'ajustement~de~Kolmogorov-Smirnov~et~Shapiro-Wilks.}}\label{faire-un-test-dajustement-de-kolmogorov-smirnov-et-shapiro-wilks.}

\subsubsection{\texorpdfstring{\textsuperscript{\ul{Peut-on~considérer~que~les~données~\emph{rivers}~sont~distribuées~selon~une~loi~normale~?}}}{Peut-on~considérer~que~les~données~rivers~sont~distribuées~selon~une~loi~normale~?}}\label{peut-on-considuxe9rer-que-les-donnuxe9es-rivers-sont-distribuuxe9es-selon-une-loi-normale}

\subsubsection{\texorpdfstring{\textsuperscript{\ul{Hypothèse~du~test~de~kolmogorov}}}{Hypothèse~du~test~de~kolmogorov}}\label{hypothuxe8se-du-test-de-kolmogorov}

Soit F la fonction de répartition de notre variable Obs1 représentatif
de rivers

\(H_o:F=F_o\)

\(H_1:F \neq F_o\)

\ul{\textbf{Code}}

\begin{Shaded}
\begin{Highlighting}[]
\NormalTok{Obs1}\OtherTok{=}\NormalTok{data2}
\CommentTok{\#test de kolmogorov simrov}
\FunctionTok{ks.test}\NormalTok{(Obs1, }\StringTok{"pnorm"}\NormalTok{, }\AttributeTok{mean =}\NormalTok{ mu1, }\AttributeTok{sd =}\NormalTok{ ecart\_type1)}
\end{Highlighting}
\end{Shaded}

\begin{verbatim}
## Warning in ks.test.default(Obs1, "pnorm", mean = mu1, sd = ecart_type1): ties
## should not be present for the one-sample Kolmogorov-Smirnov test
\end{verbatim}

\begin{verbatim}
## 
##  Asymptotic one-sample Kolmogorov-Smirnov test
## 
## data:  Obs1
## D = 0.20825, p-value = 9.768e-06
## alternative hypothesis: two-sided
\end{verbatim}

\textsuperscript{\ul{\textbf{Resultats}}}

\begin{verbatim}
Asymptotic one-sample Kolmogorov-Smirnov test  data:  Obs1 D = 0.20825, p-value = 9.768e-06 alternative hypothesis: two-sided
\end{verbatim}

\subsubsection{\texorpdfstring{\textsuperscript{\ul{Hypothèse~du~test~de~shapiro}}}{Hypothèse~du~test~de~shapiro}}\label{hypothuxe8se-du-test-de-shapiro}

\(H_o:\) la variable Obs1 suit une loi normale

\(H_1:\) la variable Obs1 ne suit pas une loi normale

\ul{\textbf{Code}}

\begin{Shaded}
\begin{Highlighting}[]
\NormalTok{Obs1}\OtherTok{=}\NormalTok{data2}
\CommentTok{\#test de shapiro}
\FunctionTok{shapiro.test}\NormalTok{(Obs1)}
\end{Highlighting}
\end{Shaded}

\begin{verbatim}
## 
##  Shapiro-Wilk normality test
## 
## data:  Obs1
## W = 0.66662, p-value < 2.2e-16
\end{verbatim}

\textsuperscript{\ul{\textbf{Resultats}}}

\begin{verbatim}
Shapiro-Wilk normality test  data:  Obs1 W = 0.66662, p-value < 2.2e-16
\end{verbatim}

Les p-value étant faible dans les deux cas en prenant \(\alpha=5\)\%
alors on rejette les hypothèse \(H_o\)

\textbf{On conclut donc qu'on ne peut pas considérer que les données
\emph{rivers} sont distribuées selon une loi normale}

\subsection{EXERCICE4}\label{exercice4}

\subsubsection{\texorpdfstring{\ul{\textbf{Répondons à la question en
supposant une hypothèse normale des
données.}}}{Répondons à la question en supposant une hypothèse normale des données.}}\label{ruxe9pondons-uxe0-la-question-en-supposant-une-hypothuxe8se-normale-des-donnuxe9es.}

Soit \(m_o\) la moyenne de l'échantillon exrait et \(m_e\) celui de
l'échantillon 24h après

\(H_o:m_o= m_e\)

\(H_1:m_o < m_e\)

Les variables x\_apres,x\_avant sont des variables appariés issus des
mêmes individus

Donc il s'agit d'un test sur la moyenne de deux echantillons de
variables aléatoires appariés

\ul{\textbf{Code}}

\begin{Shaded}
\begin{Highlighting}[]
\CommentTok{\#Données}
\NormalTok{x\_avant}\OtherTok{=}\FunctionTok{c}\NormalTok{(}\DecValTok{12000}\NormalTok{,}\DecValTok{13000}\NormalTok{,}\DecValTok{21500}\NormalTok{,}\DecValTok{17000}\NormalTok{,}\DecValTok{15000}\NormalTok{,}\DecValTok{22000}\NormalTok{,}\DecValTok{11000}\NormalTok{,}\DecValTok{21000}\NormalTok{)}

\NormalTok{x\_apres}\OtherTok{=}\FunctionTok{c}\NormalTok{(}\DecValTok{11000}\NormalTok{,}\DecValTok{20000}\NormalTok{,}\DecValTok{31000}\NormalTok{,}\DecValTok{28000}\NormalTok{,}\DecValTok{26000}\NormalTok{,}\DecValTok{30000}\NormalTok{,}\DecValTok{16000}\NormalTok{,}\DecValTok{29000}\NormalTok{)}


\CommentTok{\# Test }

\FunctionTok{t.test}\NormalTok{(x\_apres,x\_avant,}\AttributeTok{paired=}\StringTok{"TRUE"}\NormalTok{)}
\end{Highlighting}
\end{Shaded}

\begin{verbatim}
## 
##  Paired t-test
## 
## data:  x_apres and x_avant
## t = 5.2786, df = 7, p-value = 0.00115
## alternative hypothesis: true mean difference is not equal to 0
## 95 percent confidence interval:
##   4036.762 10588.238
## sample estimates:
## mean difference 
##          7312.5
\end{verbatim}

\ul{\textbf{Conclusion}}

\begin{verbatim}
p-value = 0.00115 
La p-value = 0.00115 <0.05 donc on rejette l'hypothèse intiale 
\end{verbatim}

Par conséquent l'hypothèse \(H_1\) est retenue .

Il existe donc un accroissement significatif du nombre de bactéries par
cm\textsuperscript{3} au cours du temps

\subsubsection{\texorpdfstring{~ \textbf{Répondons à la question par un
test du
signe.}}{~ Répondons à la question par un test du signe.}}\label{ruxe9pondons-uxe0-la-question-par-un-test-du-signe.}

\begin{Shaded}
\begin{Highlighting}[]
\CommentTok{\# Médiane hypothétique }

\NormalTok{mediane\_theorique }\OtherTok{\textless{}{-}} \FunctionTok{mean}\NormalTok{(x\_avant)}

\CommentTok{\# On compte les signes }

\NormalTok{plus }\OtherTok{\textless{}{-}} \FunctionTok{sum}\NormalTok{(x\_apres}\SpecialCharTok{\textgreater{}}\NormalTok{ mediane\_theorique) }

\NormalTok{moins }\OtherTok{\textless{}{-}} \FunctionTok{sum}\NormalTok{(x\_apres }\SpecialCharTok{\textless{}}\NormalTok{ mediane\_theorique)  }

\NormalTok{n }\OtherTok{\textless{}{-}}\NormalTok{ plus }\SpecialCharTok{+}\NormalTok{ moins  }\CommentTok{\# On ignore les valeurs égales à la médiane }

\CommentTok{\# Test du signe unilatéral }

\FunctionTok{binom.test}\NormalTok{(plus, }\AttributeTok{n =}\NormalTok{ n, }\AttributeTok{p =} \FloatTok{0.5}\NormalTok{, }\AttributeTok{alternative =} \StringTok{"greater"}\NormalTok{)}
\end{Highlighting}
\end{Shaded}

\begin{verbatim}
## 
##  Exact binomial test
## 
## data:  plus and n
## number of successes = 6, number of trials = 8, p-value = 0.1445
## alternative hypothesis: true probability of success is greater than 0.5
## 95 percent confidence interval:
##  0.4003106 1.0000000
## sample estimates:
## probability of success 
##                   0.75
\end{verbatim}

\paragraph{\texorpdfstring{\ul{\textbf{Conclusion:}}}{Conclusion:}}\label{conclusion}

On conclut grace a la probabilité de succès de 75\% qu'il y a
accroissement significatif du nombre de bactéries par
cm\textsuperscript{3}

au cours du temps.

\subsubsection{Répondons à la question par un test de
Mann-Whitney}\label{ruxe9pondons-uxe0-la-question-par-un-test-de-mann-whitney}

\begin{Shaded}
\begin{Highlighting}[]
\CommentTok{\# Test de Wilcoxon{-}Mann{-}Whitney (non paramétrique) }


\FunctionTok{wilcox.test}\NormalTok{(x\_apres,x\_avant,}\AttributeTok{alternative =} \StringTok{"greater"}\NormalTok{, }\AttributeTok{paired =} \ConstantTok{TRUE}\NormalTok{) }
\end{Highlighting}
\end{Shaded}

\begin{verbatim}
## Warning in wilcox.test.default(x_apres, x_avant, alternative = "greater", :
## impossible de calculer la p-value exacte avec des ex-aequos
\end{verbatim}

\begin{verbatim}
## 
##  Wilcoxon signed rank test with continuity correction
## 
## data:  x_apres and x_avant
## V = 35, p-value = 0.01028
## alternative hypothesis: true location shift is greater than 0
\end{verbatim}

\ul{\textbf{Conclusion}}

\begin{verbatim}
p-value = 0.01028

La p-value < 0.05 donc on rejette l'hypothèse Ho 
 
\end{verbatim}

Par conséquent on retient que \(H_1\) est vrai donc il existe un
accroissement significatif du nombre de bactéries par
cm\textsuperscript{3}

au cours du temps.

\end{document}
